L'operateur quasi-lineaire est caracterise par un tenseur de diffusion
\[
\left\{ \begin{array}{l}
\widetilde{D}_{pp}^{\text{RF}}=(1-\xi^{2})D_{(n)}^{\text{RF}}(p,\xi)\\
\widetilde{D}_{p\xi}^{\text{RF}}=-\frac{\sqrt{1-\xi^{2}}}{\xi}\left[1-\xi^{2}-\frac{n\Omega}{\omega}\right]D_{(n)}^{\text{RF}}(p,\xi)\\
\widetilde{D}_{\xi p}^{\text{RF}}=-\frac{\sqrt{1-\xi^{2}}}{\xi}\left[1-\xi^{2}-\frac{n\Omega}{\omega}\right]D_{(n)}^{\text{RF}}(p,\xi)\\
\widetilde{D}_{\xi\xi}^{\text{RF}}=\frac{1}{\xi^{2}}\left[1-\xi^{2}-\frac{n\Omega}{\omega}\right]^{2}D_{(n)}^{\text{RF}}(p,\xi)\end{array}\right.\]


On effectue la moyenne sur la periode de rebond \begin{eqnarray*}
\widetilde{D}_{pp}^{\text{RF}(0)} & = & (1-\xi_{0}^{2})\sigma\left\{ \frac{\sigma\xi}{\xi_{0}}D_{(n)}^{\text{RF}}(p,\xi)\right\} \\
\widetilde{D}_{p\xi}^{\text{RF}(0)} & = & -\frac{\sqrt{1-\xi_{0}^{2}}}{\xi_{0}}\left[1-\xi_{0}^{2}-\frac{n\Omega_{0}}{\omega}\right]\sigma\left\{ \frac{\sigma\xi}{\xi_{0}}D_{(n)}^{\text{RF}}(p,\xi)\right\} \\
\widetilde{F}_{p}^{\text{RF}(0)} & = & -\frac{\sqrt{1-\xi_{0}^{2}}}{p\xi_{0}^{3}}\frac{\sqrt{1-\xi_{0}^{2}}}{\xi_{0}}\left[1-\xi_{0}^{2}-\frac{n\Omega_{0}}{\omega}\right]\sigma\left\{ \left(\Psi(\theta)-1\right)\frac{\sigma\xi_{0}}{\xi}D_{(n)}^{\text{RF}}(p,\xi)\right\} \end{eqnarray*}
\begin{eqnarray*}
\widetilde{D}_{\xi p}^{\text{RF}(0)} & = & -\frac{\sqrt{1-\xi_{0}^{2}}}{\xi_{0}}\left[1-\xi_{0}^{2}-\frac{n\Omega_{0}}{\omega}\right]\sigma\left\{ \frac{\sigma\xi}{\xi_{0}}D_{(n)}^{\text{RF}}(p,\xi)\right\} \\
\widetilde{D}_{\xi\xi}^{\text{RF}(0)} & = & \frac{1}{\xi_{0}^{2}}\left[1-\xi_{0}^{2}-\frac{n\Omega_{0}}{\omega}\right]^{2}\sigma\left\{ \frac{\sigma\xi}{\xi_{0}}D_{(n)}^{\text{RF}}(p,\xi)\right\} \\
\widetilde{F}_{\xi}^{\text{RF}(0)} & = & \frac{\sqrt{1-\xi_{0}^{2}}}{p\xi_{0}^{3}}\frac{1}{\xi_{0}^{2}}\left[1-\xi_{0}^{2}-\frac{n\Omega_{0}}{\omega}\right]^{2}\sigma\left\{ \left(\Psi(\theta)-1\right)\frac{\sigma\xi_{0}}{\xi}D_{(n)}^{\text{RF}}(p,\xi)\right\} \end{eqnarray*}


Par definition \begin{eqnarray*}
\widetilde{D}_{(n)}^{\text{RF}(0)}(p,\xi_{0}) & = & \sigma\left\{ \frac{\sigma\xi}{\xi_{0}}D_{(n)}^{\text{RF}}(p,\xi)\right\} \\
\widetilde{D}_{F(n)}^{\text{RF}(0)}(p,\xi_{0}) & = & \sigma\left\{ \left(\Psi(\theta)-1\right)\frac{\sigma\xi_{0}}{\xi}D_{(n)}^{\text{RF}}(p,\xi)\right\} \end{eqnarray*}


La discretisation des elements du tenseur quasi-lineaire est alors
immediate, et les coefficients qui interviennent dans la matrice $\widetilde{\overline{\overline{M}}}^{\left(0\right)}$sont
\[
\widetilde{\overline{\overline{D}}}^{\text{RF}(0)}\rightarrow\left\{ \begin{array}{l}
\widetilde{D}_{pp,i+1,j+1/2}^{\text{RF}(0)}=(1-\xi_{0j+1/2}^{2})\widetilde{D}_{(n),i+1,j+1/2}^{\text{ RF}(0)}\\
\widetilde{D}_{pp,i,j+1/2}^{\text{RF}(0)}=(1-\xi_{0j+1/2}^{2})\widetilde{D}_{(n),i,j+1/2}^{\text{RF}(0)}\\
\widetilde{D}_{p\xi,i+3/2,j+1/2}^{\text{RF}\left(0\right)}=-\sqrt{1-\xi_{0j+1/2}^{2}}/\xi_{0j+1/2}\left[1-\xi_{0j+1/2}^{2}-n\Omega_{0}/\omega\right]\widetilde{D}_{(n),i+3/2,j+1/2}^{\text{RF}(0)}\\
\widetilde{D}_{\xi p,i+1/2,j+3/2}^{\text{RF}\left(0\right)}=-\sqrt{1-\xi_{0j+1}^{2}}/\xi_{0j+1}\left[1-\xi_{0j+1}^{2}-n\Omega_{0}/\omega\right]\widetilde{D}_{(n),i+1/2,j+3/2}^{\text{RF}(0)}\\
\widetilde{D}_{p\xi,i+1/2,j+1/2}^{\text{RF}\left(0\right)}=-\sqrt{1-\xi_{0j+1/2}^{2}}/\xi_{0j+1/2}\left[1-\xi_{0j+1/2}^{2}-n\Omega_{0}/\omega\right]\widetilde{D}_{(n),i+1/2,j+1/2}^{\text{RF}(0)}\\
\widetilde{D}_{\xi p,i+1/2,j+1/2}^{\text{RF}\left(0\right)}=-\sqrt{1-\xi_{0j+1}^{2}}/\xi_{0j+1}\left[1-\xi_{0j+1}^{2}-n\Omega_{0}/\omega\right]\widetilde{D}_{(n),i+1/2,j+1/2}^{\text{RF}(0)}\\
\widetilde{D}_{p\xi,i-1/2,j+1/2}^{\text{RF}\left(0\right)}=-\sqrt{1-\xi_{0j+1/2}^{2}}/\xi_{0j+1/2}\left[1-\xi_{0j+1/2}^{2}-n\Omega_{0}/\omega\right]\widetilde{D}_{(n),i-1/2,j+1/2}^{\text{RF}(0)}\\
\widetilde{D}_{\xi p,i+1/2,j-1/2}^{\text{RF}\left(0\right)}=-\sqrt{1-\xi_{0j+1}^{2}}/\xi_{0j+1}\left[1-\xi_{0j+1}^{2}-n\Omega_{0}/\omega\right]\widetilde{D}_{(n),i+1/2,j-1/2}^{\text{RF}(0)}\\
\widetilde{D}_{\xi\xi,i+1/2,j+1}^{\text{RF}(0)}=1/\xi_{0j+1}^{2}\left[1-\xi_{0j+1}^{2}-n\Omega_{0}/\omega\right]\widetilde{D}_{(n),i+1/2,j+1}^{\text{RF}(0)}\\
\widetilde{D}_{\xi\xi,i+1/2,j}^{\text{RF}(0)}=1/\xi_{0j}^{2}\left[1-\xi_{0j}^{2}-n\Omega_{0}/\omega\right]\widetilde{D}_{(n),i+1/2,j}^{\text{RF}(0)}\end{array}\right.\]
 et \[
\widetilde{\overline{\overline{F}}}^{\text{RF}(0)}\rightarrow\left\{ \begin{array}{l}
\widetilde{F}_{p,i+1,j+1/2}^{\text{RF}(0)}=-\left(1-\xi_{0j+1/2}^{2}\right)/\left(p_{i+1}\xi_{0j+1/2}^{4}\right)\left[1-\xi_{0j+1/2}^{2}-n\Omega_{0}/\omega\right]\widetilde{D}_{F(n),i+1,j+1/2}^{\text{ RF}(0)}\\
\widetilde{F}_{p,i,j+1/2}^{\left(0\right)}=-\left(1-\xi_{0j+1/2}^{2}\right)/\left(p_{i}\xi_{0j+1/2}^{4}\right)\left[1-\xi_{0j+1/2}^{2}-n\Omega_{0}/\omega\right]\widetilde{D}_{F(n),i,j+1/2}^{\text{ RF}(0)}\\
\widetilde{F}_{\xi,i+1/2,j+1}^{\left(0\right)}=\sqrt{1-\xi_{0j+1}^{2}}/\left(p_{i+1/2}\xi_{0j+1}^{5}\right)\left[1-\xi_{0j+1}^{2}-n\Omega_{0}/\omega\right]^{2}\widetilde{D}_{F(n),i+1/2,j+1}^{\text{RF}(0)}\\
\widetilde{F}_{\xi,i+1/2,j}^{\left(0\right)}=\sqrt{1-\xi_{0j}^{2}}/\left(p_{i+1/2}\xi_{0j}^{5}\right)\left[1-\xi_{0j}^{2}-n\Omega_{0}/\omega\right]^{2}\widetilde{D}_{F(n),i+1/2,j}^{\text{RF}(0)}\end{array}\right.\]


Sous cette forme, le calcul garde une tres grande generalite, quelque
soit le type d'onde.
